\subsubsection{Validation Curve}

Để đánh giá một cách hệ thống ảnh hưởng của độ phức tạp mô hình đến khả năng tổng quát hóa,
chúng tôi xây dựng các validation curve cho ba nhóm mô hình phi tuyến: Polynomial Basis,
Radial Basis Function (RBF) và Fourier Basis. 

Hiệu năng mô hình được đo bằng sai số bình phương trung bình (Mean Squared Error - MSE)
trên cả tập huấn luyện và tập validation, nhằm quan sát sự đánh đổi giữa bias và variance
khi độ phức tạp tăng dần.

Về mặt lý thuyết, bài toán có thể được xem như:
\begin{equation}
\hat{f} = \arg\min_f \boldsymbol{E} \left[ (y - f(x))^2 \right]
\end{equation}
trong đó sai số có thể phân rã thành:

\begin{equation}
\text{Error} = \text{Bias}^2 + \text{Variance} + \text{Noise}
\end{equation}

%------------------------------------------------
\paragraph{1. Polynomial Basis (theo bậc đa thức $p$)}

Trong thí nghiệm này, chúng tôi mở rộng không gian đặc trưng theo dạng đa thức:
\[
\phi(x) = [1, x, x^2, ..., x^p]
\]

và khảo sát ảnh hưởng của bậc $p$ lên hiệu năng mô hình.

\begin{figure}[htbp]
    \centering
    \includegraphics[width=0.75\linewidth]{poly_validation_curve.png}
    \caption{Validation curve: MSE theo bậc đa thức $p$}
    \label{fig:poly_curve}
\end{figure}

\paragraph{Phân tích chi tiết:}

Khi $p = 1$, mô hình tương đương hồi quy tuyến tính đơn giản, đạt MSE thấp khoảng $0.5349$,
cho thấy giả thiết tuyến tính ban đầu phù hợp tương đối với dữ liệu.

Tuy nhiên, khi tăng $p$, ma trận thiết kế $X^TX$ trở nên ngày càng kém điều kiện (ill-conditioned),
do sự xuất hiện của các cột $x^p$ có độ lớn tăng theo cấp số nhân:
\[
|x^p| \gg |x^{p-1}|
\]

Điều này dẫn đến hiện tượng khuếch đại sai số trong nghiệm:
\[
\hat{w} = (X^TX)^{-1}X^Ty
\]

Kết quả thực nghiệm cho thấy:
\begin{itemize}
    \item $p = 2$: MSE tăng nhẹ nhưng vẫn ổn định ($2.0967$)
    \item $p = 3$: MSE tăng đột biến ($65.1358$)
    \item $p \ge 4$: MSE bùng nổ theo cấp số mũ (lên tới $10^{12}$)
\end{itemize}

Hiện tượng này phản ánh rõ ràng sự mất ổn định số học (numerical instability)
và overfitting nghiêm trọng do mô hình cố gắng nội suy chính xác từng điểm dữ liệu.

Ngoài ra, Polynomial basis còn chịu ảnh hưởng mạnh bởi hiện tượng Runge phenomenon,
gây dao động lớn ở biên miền dữ liệu.

%------------------------------------------------
\paragraph{2. RBF Basis (theo số tâm $M$)}

Trong mô hình RBF, hàm cơ sở được định nghĩa:
\[
\phi_i(x) = \exp\left(-\frac{\|x - c_i\|^2}{2\sigma^2}\right)
\]

Trong đó $c_i$ là các tâm (centers) và $\sigma$ điều khiển độ lan tỏa.

\begin{figure}[htbp]
    \centering
    \includegraphics[width=0.75\linewidth]{rbf_validation_curve.png}
    \caption{Validation curve: MSE theo số tâm RBF}
    \label{fig:rbf_curve}
\end{figure}

\paragraph{Phân tích chi tiết:}

Khi số lượng tâm nhỏ ($M=5$), mô hình có bias cao do không gian biểu diễn còn hạn chế,
dẫn đến MSE tương đối lớn ($1.0185$).

Khi tăng $M$, không gian hàm cơ sở trở nên giàu biểu diễn hơn, giúp giảm bias:
\[
\text{Bias} \downarrow \quad \text{when } M \uparrow
\]

Đồng thời, do tính chất local support của RBF, các hàm cơ sở có độ tương quan thấp:
\[
\text{Cov}(\phi_i, \phi_j) \approx 0 \quad (i \neq j)
\]

Điều này giúp kiểm soát variance hiệu quả.

Kết quả thực nghiệm:
\begin{itemize}
    \item MSE giảm đều từ $1.0185 \rightarrow 0.3966$
    \item Không xuất hiện hiện tượng bùng nổ sai số
    \item Mô hình đạt tối ưu quanh $M = 85$--$95$
\end{itemize}

Như vậy, RBF thể hiện tính ổn định cao nhờ cấu trúc kernel cục bộ,
giúp cân bằng tốt giữa bias và variance.

%------------------------------------------------
\paragraph{3. Fourier Basis (theo số hàm cơ sở $M$)}

Fourier basis mở rộng hàm theo miền tần số:
\[
\phi_k(x) = \sin(kx), \quad \cos(kx)
\]

với $k$ đóng vai trò tần số (frequency component).

\begin{figure}[htbp]
    \centering
    \includegraphics[width=0.75\linewidth]{fourier_validation_curve.png}
    \caption{Validation curve: MSE theo số hàm cơ sở Fourier}
    \label{fig:fourier_curve}
\end{figure}

\paragraph{Phân tích chi tiết:}

Khi $M$ nhỏ, mô hình chỉ biểu diễn được các thành phần tần số thấp,
tương ứng với hàm trơn (low-frequency signal), dẫn đến underfitting:
\[
\text{MSE} \approx 0.55
\]

Khi $M$ tăng lên khoảng $3$--$5$, mô hình đạt trạng thái tối ưu:
\[
\text{MSE}_{min} \approx 0.42
\]

Tại đây, mô hình cân bằng tốt giữa:
\begin{itemize}
    \item thành phần tín hiệu (signal)
    \item và độ phức tạp biểu diễn
\end{itemize}

Tuy nhiên, khi $M > 6$, các thành phần tần số cao (high-frequency components)
bắt đầu chi phối mô hình. Điều này dẫn đến việc mô hình học cả nhiễu (noise fitting),
gây ra hiện tượng spectral overfitting.

Thực nghiệm cho thấy:
\begin{itemize}
    \item $M=8$: MSE tăng mạnh ($10.8597$)
    \item $M=9$: MSE bùng nổ ($381.1882$)
    \item $M=10$: MSE cực lớn ($3791.8363$)
\end{itemize}

Hiện tượng này tương tự Gibbs phenomenon trong phân tích Fourier,
khi khai triển chuỗi Fourier với số hạng hữu hạn trên dữ liệu có nhiễu.

%------------------------------------------------
\paragraph{Nhận xét tổng quát}

Từ ba nhóm mô hình, có thể rút ra các kết luận quan trọng:

\begin{itemize}
    \item \textbf{Polynomial Basis}: có độ nhạy cao với bậc $p$, dễ gây mất ổn định số học
    và overfitting nghiêm trọng do tính toàn cục của các hàm cơ sở.

    \item \textbf{RBF Basis}: thể hiện tính ổn định cao nhờ tính cục bộ hóa (localization),
    giúp kiểm soát tốt variance và giảm thiểu hiện tượng overfitting.

    \item \textbf{Fourier Basis}: hiệu quả tốt trong miền tần số thấp nhưng dễ bị
    overfitting khi mở rộng sang miền tần số cao do học cả nhiễu.
\end{itemize}

Tổng thể, các kết quả thực nghiệm khẳng định rằng việc tăng độ phức tạp mô hình
không đảm bảo cải thiện hiệu năng, mà cần tuân theo nguyên lý \textit{bias-variance tradeoff}
và đặc tính cấu trúc của không gian hàm cơ sở (function space geometry).